\chapter{Pain Medicine and Headache Disorders}
\label{ch:pain}

%-------------------------------------------------------------
\section{Chronic Pain Syndromes}
%-------------------------------------------------------------

%-------------------------------------------------------------

%-------------------------------------------------------------

%-------------------------------------------------------------

%-------------------------------------------------------------

%-------------------------------------------------------------

%-------------------------------------------------------------

%-------------------------------------------------------------

\label{sec:Chronic Pain Syndromes}
%-------------------------------------------------------------\label{sec:Chronic Pain Syndromes}

\subsection{Other Primary Headache Disorders}
\label{sec:Other Primary Headache Disorders}

Primary cough headache is a primary headache disorder brought on by coughing, straining, or Valsalva maneuver, typically bilateral, sudden-onset, lasting 1 second to 30 minutes, with a prevalence of approximately 1\%. Secondary causes (Chiari malformation, intracranial hypertension) must be excluded by MRI. Indomethacin is first-line treatment. Prognosis is favorable.

Primary exertional headache is a primary headache disorder precipitated by any form of exercise, typically bilateral, pulsating, lasting 5 minutes to 48 hours, more common in hot weather, at high altitude, and in patients with a history of migraine. Treatment includes ergonovine, indomethacin, or triptans taken prior to exercise. Prognosis is favorable.

Primary stabbing headache is a primary headache disorder presenting with brief, sharp, stabbing pains (``ice pick headache'') lasting 1--10 seconds, occurring in single stabs or a series, confined to the first division of the trigeminal nerve. Indomethacin is first-line. Prognosis is favorable.

Hypnic headache (``alarm clock headache'') is a rare primary headache disorder in patients over 50 years of age, with attacks strictly sleep-related, awakening the patient at a consistent time (typically 1--3 hours after sleep onset), bilateral, dull, lasting 15--180 minutes, without autonomic features. Caffeine, lithium, or indomethacin are used. Prognosis is generally favorable.

\subsection{Paroxysmal Hemicrania}
\label{sec:Paroxysmal Hemicrania}
Paroxysmal hemicrania is a trigeminal autonomic cephalalgia characterized by strictly unilateral, severe pain attacks lasting 2--30 minutes, occurring several to many times per day (typically >5/day). Onset is typically in adulthood, with a slight female predominance. The condition results from activation of the trigeminovascular system with hypothalamic dysfunction. Clinical features include unilateral pain with ipsilateral autonomic features identical to cluster headache. Diagnosis requires a therapeutic trial of indomethacin; the defining feature is absolute response to indomethacin (75--150 mg/day). Management involves indomethacin with gastrointestinal protection. Prognosis is favorable with treatment; episodic and chronic forms exist.

\subsection{Phantom Limb Pain}
\label{sec:Phantom Limb Pain}
Phantom limb pain is pain perceived in a body part that has been amputated. It occurs in 60--80\% of amputees, typically within the first week after amputation, and persists in 50--80\% at 2 years. The condition results from peripheral neuroma formation with ectopic discharges, spinal cord reorganization, and cortical reorganization (invasion of the deafferented somatosensory cortex by adjacent cortical inputs). Clinical features include pain perceived in the absent limb, distinct from residual limb pain and non-painful phantom limb sensations. Diagnosis is clinical. Management is challenging; first-line treatments include gabapentin, pregabalin, TCAs, and tramadol; non-pharmacological approaches include mirror therapy, graded motor imagery, and virtual reality therapy. Prognosis is variable; treatment is often challenging.

\subsection{Physical and Rehabilitation Approaches}
\label{sec:Physical and Rehabilitation Approaches}

Physical therapy is essential for chronic pain management. Therapeutic exercise includes range of motion, stretching, strengthening, and aerobic conditioning. Graded activity programs gradually increase activity levels based on time rather than pain, countering fear-avoidance behavior. Manual therapy (joint mobilization, soft tissue mobilization, myofascial release) addresses musculoskeletal dysfunction. Modalities such as transcutaneous electrical nerve stimulation (TENS), ultrasound, heat, cold, and traction provide short-term symptomatic relief. Occupational therapy focuses on activity modification, ergonomics, energy conservation, and adaptive equipment to improve function in daily activities.

\subsection{Post-Traumatic Headache}
\label{sec:Post-Traumatic Headache}
Post-traumatic headache (PTH) is a secondary headache disorder and the most common symptom following traumatic brain injury (TBI), affecting 30--90\% of patients. It is classified as acute (<3 months) or persistent ($\ge$3 months). The condition results from a neurometabolic cascade after TBI including axonal injury, excitotoxicity, mitochondrial dysfunction, altered cerebral blood flow, and activation of the trigeminovascular system. Clinical features most commonly mimic migraine or tension-type headache. Diagnosis is clinical with history of TBI. Management follows standard primary headache guidelines with amitriptyline, topiramate, and gabapentin commonly used; physical therapy for associated cervical dysfunction is essential. Prognosis is variable; most improve within 3--6 months but some remain symptomatic for years.

\subsection{Postherpetic Neuralgia}
\label{sec:Postherpetic Neuralgia}
Postherpetic neuralgia (PHN) is neuropathic pain persisting for $\ge$3 months after the resolution of herpes zoster (shingles) rash. It occurs in 10--20\% of patients with herpes zoster, increasing to >30\% in patients over 60 years of age. The condition results from peripheral and central sensitization from varicella-zoster virus reactivation in the dorsal root ganglia, leading to nerve damage, inflammation, and neuronal hyperexcitability. Clinical features include burning, aching, electric, and lancinating pain with allodynia and hyperalgesia in a dermatomal distribution (most commonly thoracic). Diagnosis is clinical. First-line pharmacotherapy includes topical lidocaine 5\% patches, pregabalin, gabapentin, and TCAs; second-line includes tramadol, capsaicin 8\% patches, and opioid analgesics. Prognosis is variable; prevention is the best treatment, with herpes zoster vaccination (recombinant zoster vaccine) reducing the incidence of PHN by approximately 90\%.

%-------------------------------------------------------------
\section{Headache Disorders --- Primary Headaches}
%-------------------------------------------------------------

%-------------------------------------------------------------

%-------------------------------------------------------------

%-------------------------------------------------------------

%-------------------------------------------------------------

%-------------------------------------------------------------

%-------------------------------------------------------------

%-------------------------------------------------------------

\label{sec:Headache Disorders --- Primary Headaches}
%-------------------------------------------------------------\label{sec:Headache Disorders Primary Headaches}

Headache disorders are among the most common disorders of the nervous system. Primary headaches are those in which the headache itself is the disorder, not a symptom of an underlying disease.

\subsection{Cervicogenic Headache}
\label{sec:Cervicogenic Headache}
Cervicogenic headache is a secondary headache originating from bony or soft tissue structures of the cervical spine. The condition results from convergence of cervical afferents (C1--C3) with trigeminal afferents in the trigeminocervical complex. Clinical features include unilateral pain starting in the neck or occiput and spreading to the frontotemporal region, precipitated or worsened by specific neck movements or sustained head postures, with reduced range of motion and tenderness over upper cervical facets or occipital nerve. Diagnosis is supported by temporary relief with anesthetic block of the C2--3 nerve root or greater occipital nerve. Management includes physical therapy, occipital nerve block, radiofrequency ablation, and optimization of ergonomics. Prognosis is variable and depends on the underlying pathology.

\subsection{Chronic Low Back Pain}
\label{sec:Chronic Low Back Pain}
Chronic low back pain (cLBP) is pain in the lower back persisting for $\ge$3 months. It is the leading cause of years lived with disability worldwide, with a lifetime prevalence of 60--80\%. Most cases are non-specific (no identifiable structural pathology). The condition results from a biopsychosocial interaction of biological factors (structural pathology, inflammation), psychological factors (catastrophizing, fear-avoidance beliefs, depression), and social factors (occupational demands, disability compensation). Clinical features include pain in the lower back that may radiate to the legs (sciatica), with or without neurological deficits. Diagnosis is clinical; red flags for serious pathology include age >50 or <20, recent trauma, unexplained weight loss, fever, night pain, history of cancer, and progressive neurological deficits. Initial management includes education, activity modification, NSAIDs, and muscle relaxants; first-line pharmacotherapy is NSAIDs, with TCAs, SNRIs, and gabapentinoids as adjuncts; non-pharmacological treatments include physical therapy, manual therapy, acupuncture, cognitive behavioral therapy, and graded activity programs; interventional procedures include epidural steroid injections, facet joint injections, and spinal cord stimulation. Prognosis is variable; most cases improve with conservative management.

\subsection{Chronic Neck Pain}
\label{sec:Chronic Neck Pain}
Chronic neck pain has a point prevalence of 10--15\% and a lifetime prevalence of 50--70\%. It is more common in women and office workers. The condition results from mechanical neck pain (poor posture, prolonged computer use), cervical spondylosis, cervical radiculopathy, or myofascial pain. Clinical features include pain in the neck that may radiate to the shoulders or arms. Diagnosis is clinical. Management parallels chronic low back pain: physical therapy, manual therapy, exercise, NSAIDs, and muscle relaxants; cervical epidural steroid injections and selective nerve root blocks are used for radicular pain. Prognosis is variable.

\subsection{Chronic Pain vs Acute Pain}
\label{sec:Chronic Pain vs Acute Pain}

Acute pain is a direct physiological response to a noxious stimulus or injury, with a reversible time course proportional to tissue healing (typically days to weeks). It serves an adaptive, protective function and resolves with treatment of the underlying cause. Chronic pain is pain persisting beyond normal tissue healing time (typically >3 months) or pain associated with a chronic condition. It is maladaptive, often lacks a clear protective function, and involves neuroplastic changes in the peripheral and central nervous systems. Chronic pain affects approximately 20--30\% of the adult population worldwide. Management requires a biopsychosocial approach rather than a purely biomedical one.

\subsection{Chronic Pelvic Pain}
\label{sec:Chronic Pelvic Pain}
Chronic pelvic pain (CPP) is persistent or recurrent pain in the pelvic region lasting $\ge$6 months in women (and $\ge$3--6 months in men). In women, estimated prevalence is 15--25\%. The condition results from various causes: in women, endometriosis, interstitial cystitis, chronic pelvic inflammatory disease, pelvic adhesions, adenomyosis, and pelvic floor myofascial pain; in men, chronic prostatitis and pudendal neuralgia. Clinical features include persistent or recurrent pelvic pain. Diagnosis requires a multidisciplinary approach including gynecological, urological, gastroenterological, and pain medicine evaluation. Management is directed at the identified cause; pelvic floor physical therapy is highly effective for pelvic floor myofascial dysfunction; pharmacotherapy includes NSAIDs, gabapentinoids, TCAs, and hormonal therapy for endometriosis. Prognosis is variable.

\subsection{Cluster Headache}
\label{sec:Cluster Headache}
Cluster headache is a trigeminal autonomic cephalalgia (TAC) characterized by strictly unilateral, severe to very severe pain in the orbital, supraorbital, temporal, or any combination of these areas, lasting 15--180 minutes if untreated, with attacks occurring from once every other day to 8 times per day. Prevalence is approximately 0.1\%, with a male predominance (3:1). The condition results from activation of the posterior hypothalamic gray matter and the trigeminovascular system. Clinical features include unilateral burning, boring, or stabbing pain with ipsilateral autonomic features (conjunctival injection, lacrimation, nasal congestion, rhinorrhea, miosis, ptosis, eyelid edema) and restlessness or agitation during attacks. Episodic cluster headache occurs in bouts with remission periods; chronic cluster headache has attacks for >1 year without remission. Diagnosis is clinical based on ICHD-3 criteria. Acute treatment includes subcutaneous sumatriptan, intranasal zolmitriptan, and high-flow oxygen. Transitional therapy uses short-course oral corticosteroids or greater occipital nerve block. Preventive therapy includes verapamil, lithium, topiramate, and valproic acid. Prognosis is chronic but attacks can be prevented during bouts with appropriate therapy.

\subsection{Complex Regional Pain Syndrome}
\label{sec:Complex Regional Pain Syndrome}
Complex regional pain syndrome (CRPS) is a painful condition characterized by disproportionate pain, sensory changes, vasomotor dysfunction, sudomotor dysfunction, and trophic changes in a limb. CRPS Type I occurs without a definable nerve lesion, typically after trauma or immobilization; CRPS Type II occurs with a definable major nerve injury. The incidence is 5--26 per 100,000 person-years, with a female-to-male ratio of 3--4:1. The condition results from multifactorial pathophysiology including peripheral and central sensitization, sympathetic nervous system dysregulation, neurogenic inflammation, and altered somatosensory processing. Clinical features are organized into four domains: sensory (allodynia, hyperalgesia), vasomotor (temperature asymmetry, skin color changes), sudomotor/edema, and motor/trophic (decreased range of motion, weakness, tremor, dystonia). Diagnosis is clinical using the Budapest diagnostic criteria. Management is most effective when initiated early and includes a multidisciplinary approach with physical therapy, pharmacotherapy (gabapentinoids, TCAs, bisphosphonates), and interventional procedures (sympathetic nerve blocks, spinal cord stimulation). Prognosis is variable; early treatment improves outcomes but up to 20\% of patients have persistent disability.

\subsection{Diabetic Peripheral Neuropathic Pain}
\label{sec:Diabetic Peripheral Neuropathic Pain}
Diabetic peripheral neuropathy (DPN) is a neuropathic pain condition affecting approximately 50\% of patients with diabetes mellitus, with 15--25\% of these patients experiencing neuropathic pain. The condition results from hyperglycemia-induced metabolic pathways (polyol pathway flux, advanced glycation end products, oxidative stress, protein kinase C activation) leading to axonal degeneration, demyelination, and loss of small and large nerve fibers. Clinical features include bilateral, symmetric, stocking-glove distribution of burning, tingling, shooting, and electric sensations, often worse at night, with reduced sensation, absent ankle reflexes, and allodynia. Diagnosis is clinical. First-line pharmacotherapy includes duloxetine, pregabalin, gabapentin, and TCAs. Prognosis is chronic; pain management does not reverse the underlying neuropathy but improves quality of life; glycemic control is essential for prevention and slowing progression.

%-------------------------------------------------------------
\section{Headache Disorders --- Secondary Headaches}
%-------------------------------------------------------------

%-------------------------------------------------------------

%-------------------------------------------------------------

%-------------------------------------------------------------

%-------------------------------------------------------------

%-------------------------------------------------------------

%-------------------------------------------------------------

%-------------------------------------------------------------

\label{sec:Headache Disorders --- Secondary Headaches}
%-------------------------------------------------------------\label{sec:Headache Disorders Secondary Headaches}

Secondary headaches are caused by an underlying disorder. They require prompt recognition as many causes are life-threatening.

\subsection{Fibromyalgia}
\label{sec:Fibromyalgia Fibromyalgia}

See Chapter \ref{ch:muscles}, Section \ref{sec:Fibromyalgia}. From a pain medicine perspective, fibromyalgia is the prototypical nociplastic pain condition involving central sensitization, dysfunctional descending pain modulation, and aberrant connectivity within the default mode network and salience network. Management emphasizes a multimodal approach: pharmacotherapy (duloxetine 60--120 mg/day, milnacipran 100--200 mg/day, pregabalin 150--600 mg/day---the three FDA-approved agents), graded aerobic exercise, cognitive behavioral therapy, sleep hygiene, and patient education. Opioids are ineffective and should be avoided.

\subsection{Glossopharyngeal Neuralgia}
\label{sec:Glossopharyngeal Neuralgia}
Glossopharyngeal neuralgia is a rare neuropathic pain disorder characterized by paroxysmal, severe, stabbing pain in the posterior pharynx, tonsillar fossa, base of tongue, and deep ear (cranial nerves IX and X). It may be caused by neurovascular compression of the glossopharyngeal nerve. Clinical features include attacks triggered by swallowing, coughing, talking, or yawning, which can be associated with bradycardia, syncope, or asystole due to vagal activation. Diagnosis is clinical. Management includes carbamazepine or oxcarbazepine as first-line; surgical options include microvascular decompression or percutaneous radiofrequency rhizotomy. Prognosis is favorable with treatment.

\subsection{Headache Attributed to Giant Cell Arteritis}
\label{sec:Headache Giant Cell Arteritis}

See Chapter \ref{ch:connective_tissue}, Section \ref{sec:Giant Cell Arteritis}.

\subsection{Headache Attributed to Intracranial Hypertension}
\label{sec:Headache Intracranial Hypertension}
Idiopathic intracranial hypertension (IIH, pseudotumor cerebri) is a secondary headache disorder characterized by elevated intracranial pressure without identifiable cause. The annual incidence is 0.5--3 per 100,000, rising to 12--20 per 100,000 in obese women of childbearing age. The condition results from impaired CSF absorption at the arachnoid granulations or elevated cerebral venous sinus pressure. Clinical features include daily, diffuse, pressure-like headache worse in the morning or with Valsalva, transient visual obscurations, pulsatile tinnitus, and papilledema; visual loss is the most serious complication. Diagnosis requires elevated opening pressure on lumbar puncture (>25 cm H2O) with normal CSF composition, normal neuroimaging, and no alternative cause. Management includes weight loss, acetazolamide, and serial ophthalmological monitoring; refractory cases may require optic nerve sheath fenestration or CSF shunting. Prognosis is variable; visual loss occurs in 10--25\% of patients.

\subsection{Headache Attributed to Intracranial Hypotension}
\label{sec:Headache Intracranial Hypotension}
Spontaneous intracranial hypotension (SIH) is a secondary headache disorder caused by spontaneous CSF leak, typically at the spinal level (thoracic or cervical). The condition results from CSF leak causing low CSF volume and pressure. Clinical features include orthostatic headache worse within seconds to minutes of upright posture and improving with recumbency, with associated neck pain, nausea, tinnitus, and dizziness. Diagnosis is by clinical presentation and MRI showing diffuse pachymeningeal enhancement, brain sagging, and subdural fluid collections. Management is conservative initially (bed rest, caffeine); epidural blood patch is the definitive treatment for persistent symptoms, with success rates of 80--90\%. Prognosis is favorable with treatment.

Post-dural puncture headache (PDPH) is a secondary headache disorder complicating lumbar puncture, spinal anesthesia, or epidural anesthesia, occurring in 1--30\% of cases. The condition results from persistent CSF leak through the dural puncture site. Clinical features include orthostatic headache typically fronto-occipital, with onset within 24--72 hours after dural puncture. Diagnosis is clinical. Management includes conservative measures for mild cases and epidural blood patch for moderate to severe cases. Prognosis is favorable.

\subsection{Headache Attributed to Sinusitis}
\label{sec:Headache Sinusitis}

See Chapter \ref{ch:nose}, Section \ref{sec:Chronic Rhinosinusitis}.

\subsection{Headache in Brain Tumors}
\label{sec:Headache Brain Tumors}

See Chapter \ref{ch:brain}, Section \ref{sec:Brain Tumors}.

\subsection{Hemicrania Continua}
\label{sec:Hemicrania Continua}
Hemicrania continua is a trigeminal autonomic cephalalgia presenting as a continuous, strictly unilateral headache of moderate intensity, with superimposed exacerbations of severe pain, present for >3 months without remission. Prevalence is approximately 0.1\% of headache clinic populations, with a female predominance. Clinical features include continuous unilateral headache with ipsilateral cranial autonomic features during exacerbations and a sense of restlessness or agitation. Diagnosis is clinical; it is exquisitely responsive to indomethacin (75--150 mg/day), with continuous pain resolving within 24--72 hours of starting therapy. Management involves indomethacin with gastrointestinal and renal monitoring. Prognosis is favorable with treatment.

\subsection{Interventional Pain Management}
\label{sec:Interventional Pain Management}

Epidural steroid injections (ESI) deliver corticosteroids into the epidural space to reduce inflammation and provide pain relief for radicular pain from disc herniation, spinal stenosis, or nerve root compression. Transforminal ESI targets specific nerve roots; interlaminar ESI provides a more diffuse effect. Complications include infection, hematoma, dural puncture, nerve injury, and intravascular injection. Fluoroscopic guidance is standard.

Peripheral nerve blocks involve injection of local anesthetic (with or without corticosteroid) around a peripheral nerve to interrupt nociceptive transmission. Common blocks include greater occipital nerve (headache), suprascapular nerve (shoulder pain), pudendal nerve (pelvic pain), and intercostal nerve (chest wall pain).

Radiofrequency ablation (RFA) applies heat generated by radiofrequency energy to produce a thermocoagulation lesion on a specific nerve, providing sustained pain relief (3--12 months). Applications include facet joint denervation (medial branch RFA for chronic axial low back pain and neck pain), sacroiliac joint denervation (lateral branch RFA), genicular nerve RFA for knee osteoarthritis, and pulsed RFA for occipital neuralgia.

Spinal cord stimulation (SCS) involves implanting electrodes in the epidural space to deliver electrical stimulation to the dorsal columns, modulating pain transmission via the gate control theory. Indications include failed back surgery syndrome, complex regional pain syndrome, painful diabetic neuropathy, and refractory angina. Modern SCS uses paresthesia-based (traditional) or paresthesia-free (high-frequency, burst) stimulation. Candidates undergo a trial period (5--7 days) before permanent implantation.

Intrathecal drug delivery (implanted pump) delivers medications directly into the cerebrospinal fluid, bypassing the blood-brain barrier and achieving high analgesic concentrations with lower systemic doses. Indications include cancer pain, severe spasticity (baclofen), and chronic non-cancer pain refractory to other treatments. Medications used include morphine, hydromorphone, ziconotide, clonidine, bupivacaine, and baclofen (for spasticity).

%-------------------------------------------------------------
\section{Neuropathic Pain Syndromes}
%-------------------------------------------------------------

%-------------------------------------------------------------

%-------------------------------------------------------------

%-------------------------------------------------------------

%-------------------------------------------------------------

%-------------------------------------------------------------

%-------------------------------------------------------------

%-------------------------------------------------------------

\label{sec:Neuropathic Pain Syndromes}
%-------------------------------------------------------------\label{sec:Neuropathic Pain Syndromes}

\subsection{Medication Overuse Headache}
\label{sec:Medication Overuse Headache}
Medication overuse headache (MOH), also called rebound headache, is a secondary headache disorder resulting from frequent or excessive use of acute headache medications. It is the most common cause of chronic daily headache, affecting 1--2\% of the general population. The condition results from frequent use of acute headache medications, particularly combination analgesics, triptans, opioids, and barbiturates. Clinical features include headache on $\ge$15 days per month with a causal relationship to medication overuse. Diagnosis requires ICHD-3 criteria. Management involves withdrawal of the overused medication (abruptly for triptans and simple analgesics, gradually for opioids and barbiturates), bridging therapy during withdrawal, and concurrent initiation of preventive therapy. Prognosis is favorable with treatment; relapse rates are 30--40\% within the first year.

\subsection{Migraine}
\label{sec:Migraine}

Migraine is a primary headache disorder affecting approximately 15\% of the global population (1 billion people), with a female-to-male ratio of 3:1. It is the second leading cause of years lived with disability worldwide. The condition results from cortical spreading depression (CSD)---a wave of neuronal depolarization followed by prolonged suppression of cortical activity---triggering the trigeminovascular system with release of calcitonin gene-related peptide (CGRP), causing vasodilation and neurogenic inflammation; genetic factors play a role, with familial hemiplegic migraine linked to mutations in CACNA1A, ATP1A2, and SCN1A. Clinical features include recurrent, unilateral (60\%), pulsating headache of moderate to severe intensity, lasting 4--72 hours, associated with nausea, vomiting, photophobia, and phonophobia; migraine with aura involves transient focal neurological symptoms preceding or accompanying the headache; chronic migraine is defined by headache on $\ge$15 days per month for >3 months. Diagnosis is clinical based on ICHD-3 criteria. Acute management involves NSAIDs or triptans for mild to moderate attacks and triptans, lasmiditan, or gepants for moderate to severe attacks. Preventive therapy is indicated for $\ge$4 migraine days per month and includes beta-blockers, amitriptyline, topiramate, candesartan, onabotulinumtoxin A for chronic migraine, and CGRP monoclonal antibodies. Prognosis is variable; migraine is a chronic condition with episodic attacks that can be managed but not cured.

\subsection{Myofascial Pain Syndrome}
\label{sec:Myofascial Pain Syndrome}
Myofascial pain syndrome (MPS) is a chronic pain syndrome characterized by the presence of trigger points---hyperirritable, taut bands of skeletal muscle that are painful on palpation and produce referred pain. The condition results from dysfunctional motor endplates with excessive acetylcholine release, sustained sarcomere contraction, local ischemia, and release of inflammatory mediators. Clinical features include regional pain (not widespread like fibromyalgia), with active trigger points causing spontaneous pain and latent trigger points painful only with palpation; common sites include the trapezius, levator scapulae, masseter, and gluteal muscles. Diagnosis is clinical. Management includes trigger point injections (dry needling or with local anesthetic), spray and stretch, physical therapy, posture correction, and ergonomic modifications; botulinum toxin injections may be used for refractory cases. Prognosis is favorable with treatment.

\subsection{New Daily Persistent Headache}
\label{sec:New Daily Persistent Headache}
New daily persistent headache (NDPH) is a primary headache disorder characterized by a sudden-onset headache that becomes continuous within 24 hours and persists for >3 months. The condition results from unclear pathophysiology, with triggers including a viral or flu-like illness (40--50\%), stressful life event, surgery, or travel. Clinical features include bilateral, pressing/tightening, moderate intensity pain; the onset is clearly remembered by the patient (``a thunderclap headache that never went away''). Diagnosis requires rigorous exclusion of secondary causes (subarachnoid hemorrhage, cerebral venous sinus thrombosis, meningitis, intracranial hypotension). Management is challenging, involving a multidisciplinary approach with amitriptyline, gabapentin, topiramate, and physical therapy. Prognosis is variable; the self-limiting subtype typically resolves within 2 years, while the refractory form causes significant disability.

\subsection{Nociception}
\label{sec:Nociception}

Nociception is the neural process of encoding noxious stimuli and involves four phases:
\begin{itemize}[noitemsep]
  \item \textbf{Transduction:} Noxious stimuli (mechanical, thermal, chemical) are converted into electrical signals at peripheral nociceptors (free nerve endings of A-delta and C fibers). Ion channels including TRPV1 (heat, capsaicin), TRPM8 (cold, menthol), TRPA1 (irritants), ASIC (acid), P2X (ATP), and voltage-gated sodium channels (NaV1.7, NaV1.8) mediate transduction.
  \item \textbf{Transmission:} The electrical signal propagates along primary afferent neurons to the dorsal horn of the spinal cord. Fast, sharp pain is carried by thinly myelinated A-delta fibers (5--30 m/s). Slow, dull, burning pain is carried by unmyelinated C fibers (0.5--2 m/s). Second-order neurons cross the midline and ascend via the spinothalamic tract to the thalamus. Third-order neurons project to the somatosensory cortex (sensory-discriminative component) and limbic system (affective-motivational component).
  \item \textbf{Modulation:} Descending pathways from the periaqueductal gray (PAG), rostral ventromedial medulla (RVM), and locus coeruleus modulate nociceptive transmission at the dorsal horn via endogenous opioids (endorphins, enkephalins, dynorphins), serotonin, and norepinephrine. This system is the basis for the analgesic effects of opioids, TCAs, and SNRIs.
  \item \textbf{Perception:} The conscious experience of pain, influenced by attention, emotion, past experience, and context. This explains the variability in pain experience across individuals and situations.
\end{itemize}

\subsection{Non-Opioid Analgesics}
\label{sec:Non-Opioid Analgesics}

Acetaminophen (paracetamol) is a centrally acting analgesic with weak anti-inflammatory effects. Its mechanism involves inhibition of cyclooxygenase (COX) in the central nervous system and interaction with the serotonergic system. It is effective for mild to moderate pain, particularly musculoskeletal pain and osteoarthritis. Maximum daily dose is 4,000 mg (3,000 mg in patients with hepatic impairment, older adults, or chronic alcohol use). Hepatotoxicity occurs at doses exceeding 4,000 mg/day or with acute overdose.

NSAIDs include non-selective COX inhibitors (ibuprofen, naproxen, diclofenac, ketorolac, indomethacin) and selective COX-2 inhibitors (celecoxib, etoricoxib). They are effective for nociceptive pain with an inflammatory component. Mechanism is inhibition of COX-1 and COX-2, reducing prostaglandin synthesis. Adverse effects include gastrointestinal ulceration and bleeding (especially with COX-1 inhibition), renal impairment (reduced prostaglandin-mediated renal blood flow), cardiovascular risk (increased in both COX-2 selective and non-selective NSAIDs), and platelet inhibition (non-selective NSAIDs). GI protection with proton pump inhibitors is recommended for patients on chronic NSAID therapy or with risk factors.

\subsection{Occipital Neuralgia}
\label{sec:Occipital Neuralgia}
Occipital neuralgia is a neuropathic pain disorder characterized by paroxysmal, stabbing, or shooting pain in the distribution of the greater, lesser, or third occipital nerves (posterior scalp, from occiput to vertex). It may be caused by entrapment of the occipital nerves by cervical muscles or by trauma, cervical arthritis, or C1--C2 pathology. Clinical features include pain accompanied by tenderness over the occipital nerve and paresthesias or dysesthesia in the affected distribution. Diagnosis is supported by temporary relief with local anesthetic block of the occipital nerve. Management includes nerve blocks with local anesthetic and corticosteroid (therapeutic and diagnostic), physical therapy, and medications (gabapentin, pregabalin, TCAs); for refractory cases, pulsed radiofrequency ablation or occipital nerve stimulation may be considered. Prognosis is variable.

\subsection{Opioid Analgesics}
\label{sec:Opioid Analgesics}

Opioids are agonists at mu, kappa, and delta opioid receptors in the central and peripheral nervous systems. They produce analgesia, euphoria, sedation, constipation, respiratory depression, nausea, and miosis. Morphine is the prototypic mu-opioid receptor agonist, effective for moderate to severe pain. It undergoes hepatic metabolism to morphine-3-glucuronide (M3G, inactive for analgesia but neuroexcitatory) and morphine-6-glucuronide (M6G, active, more potent than morphine). Hydromorphone is a semi-synthetic opioid 5--10 times more potent than morphine, with similar efficacy and safety profile. Fentanyl is a highly potent synthetic opioid (100 times more potent than morphine) with rapid onset and short duration of action, available in IV, transdermal, intranasal, and transmucosal formulations. The transdermal fentanyl patch provides continuous delivery for 72 hours and is not recommended for opioid-na{\"i}ve patients. Oxycodone is a semi-synthetic opioid available as immediate-release and extended-release formulations, often combined with acetaminophen or naloxone. Methadone is a synthetic opioid with mu-receptor agonism and NMDA receptor antagonism. It has a long and variable half-life (15--60 hours) requiring cautious titration due to risk of accumulation and QT prolongation (ECG monitoring recommended). Tapentadol is a mu-opioid receptor agonist and norepinephrine reuptake inhibitor, effective for both nociceptive and neuropathic pain. Buprenorphine is a partial mu-opioid receptor agonist and kappa-receptor antagonist, with a ceiling effect for respiratory depression. It is available in sublingual, buccal, transdermal, and implantable formulations. It may be used for chronic pain (especially in patients with opioid use disorder) and has a favorable safety profile.

\subsection{Opioid Therapy Risks}
\label{sec:Opioid Therapy Risks}

Tolerance is the reduced analgesic effect with repeated use, requiring dose escalation to maintain pain relief. It results from mu-receptor desensitization, upregulation of adenylyl cyclase, and NMDA receptor activation. Tolerance develops more slowly for analgesia than for respiratory depression, euphoria, and sedation.

Physical dependence is a state of adaptation manifested by a withdrawal syndrome when the opioid is discontinued, the dose is reduced, or an opioid antagonist is administered. Withdrawal symptoms are unpleasant but not life-threatening (except in neonates) and include anxiety, irritability, mydriasis, lacrimation, rhinorrhea, piloerection, diaphoresis, myalgias, arthralgias, abdominal cramping, diarrhea, nausea, and vomiting. Withdrawal is managed by gradual dose tapering.

Addiction (opioid use disorder, OUD) is a chronic, relapsing neurobiological disease characterized by impaired control over opioid use, compulsive use, continued use despite harm, and craving. The prevalence of OUD among patients prescribed opioids for chronic pain is estimated at 8--12\%. Risk factors include personal or family history of substance use disorder, younger age, untreated psychiatric conditions, and history of childhood abuse or trauma. Monitoring tools include prescription drug monitoring programs, urine drug testing, and risk assessment instruments (Opioid Risk Tool, Screener and Opioid Assessment for Patients with Pain). Risk mitigation strategies include avoiding high-dose therapy (morphine milligram equivalents >90 mg/day is associated with increased risk of overdose), avoiding concurrent benzodiazepine use, providing naloxone education and prescription, and establishing treatment agreements.

Opioid-induced hyperalgesia (OIH) is a paradoxical increase in pain sensitivity caused by opioid exposure, distinct from tolerance. It presents as worsening pain despite increasing opioid dose, with diffuse pain that extends beyond the original pain site. The mechanism involves activation of central glutamatergic systems (NMDA receptors), descending facilitation from the RVM, and increased dynorphin and substance P in the spinal cord. Management includes opioid dose reduction, rotation to a different opioid (methadone or buprenorphine), addition of NMDA receptor antagonists (ketamine, methadone), and non-opioid analgesic optimization.

%-------------------------------------------------------------
\section{Pain Management Principles}
%-------------------------------------------------------------

%-------------------------------------------------------------

%-------------------------------------------------------------

%-------------------------------------------------------------

%-------------------------------------------------------------

%-------------------------------------------------------------

%-------------------------------------------------------------

%-------------------------------------------------------------

\label{sec:Pain Management Principles}
%-------------------------------------------------------------\label{sec:Pain Management Principles}

\subsection{Psychological Approaches}
\label{sec:Psychological Approaches}

Cognitive behavioral therapy (CBT) is the most evidence-based psychological treatment for chronic pain. It aims to identify and modify maladaptive thoughts (catastrophizing, fear-avoidance beliefs) and behaviors (activity avoidance, pain behaviors) through cognitive restructuring, behavioral activation, graded exposure, and relaxation techniques. The goal is improved function and quality of life, not necessarily pain elimination. Biofeedback uses physiological monitoring (EMG, heart rate variability, skin conductance, breathing) to teach patients voluntary control of physiological processes that contribute to pain and stress. Mindfulness-based stress reduction (MBSR) and acceptance and commitment therapy (ACT) cultivate present-moment awareness, acceptance of pain, and value-driven behavior. These approaches reduce pain-related distress and improve function and are recommended as part of multimodal chronic pain treatment.

\subsection{SUNCT and SUNA}
\label{sec:SUNCT SUNA}
Short-lasting Unilateral Neuralgiform headache attacks with Conjunctival Injection and Tearing (SUNCT) and Short-lasting Unilateral Neuralgiform headache attacks with cranial autonomic symptoms (SUNA) are trigeminal autonomic cephalalgias characterized by brief, stabbing, burning, or electric shock-like pain in the orbital, supraorbital, or temporal region. They are more common in men (1.5--2:1), with onset typically after age 50. Clinical features include attacks lasting 1--600 seconds (typically 5--240 seconds) occurring 1--200 times per day. Diagnosis is clinical. Management includes lamotrigine as first-line for prevention, with gabapentin, topiramate, and carbamazepine as alternatives; surgical options are reserved for refractory cases. Prognosis is variable; the condition is often refractory to treatment.

\subsection{Temporomandibular Joint Disorders}
\label{sec:Temporomandibular Joint Disorders}
Temporomandibular joint disorders (TMD/TMJ) encompass a group of conditions affecting the temporomandibular joint, masticatory muscles, and associated structures. Prevalence is 5--12\%, with a female predominance (4:1) and peak incidence in reproductive years. The condition results from multifactorial causes including occlusal factors, parafunctional habits, trauma, psychosocial stress, and genetic predisposition. Clinical features include jaw pain, clicking or crepitus, limited mouth opening, locking, headache, and ear pain. Diagnosis is clinical with palpation of TMJ and muscles and range of motion assessment; imaging is reserved for atypical presentations. Conservative management is first-line: patient education, self-management, physical therapy, occlusal splints, and NSAIDs; for persistent pain, TCAs, muscle relaxants, trigger point injections, or botulinum toxin may be used. Prognosis is favorable with conservative management.

\subsection{Tension-Type Headache}
\label{sec:Tension-Type Headache}
Tension-type headache (TTH) is a primary headache disorder and the most prevalent headache disorder, with a lifetime prevalence of 30--78\%. It is more common in women and peaks in the 30--40 age range. The condition results from peripheral myofascial pain mechanisms in episodic TTH and central sensitization in chronic TTH, with muscle tenderness (pericranial, cervical, shoulder girdle) as a prominent feature. Clinical features include bilateral, pressing/tightening, non-pulsating headache of mild to moderate intensity, lasting 30 minutes to 7 days, not aggravated by routine physical activity, and not associated with nausea or vomiting. Diagnosis is clinical based on ICHD-3 criteria. Acute management includes simple analgesics (ibuprofen, naproxen, acetaminophen, aspirin). Preventive therapy for chronic TTH includes amitriptyline, mirtazapine, or venlafaxine; non-pharmacological approaches include physical therapy, relaxation training, biofeedback, and cognitive behavioral therapy. Prognosis is favorable; episodic TTH typically resolves with treatment.

\subsection{Trigeminal Autonomic Cephalalgias}
\label{sec:Trigeminal Autonomic Cephalalgias}

The trigeminal autonomic cephalalgias (TACs) are a group of primary headache disorders characterized by unilateral trigeminal distribution of pain with ipsilateral cranial autonomic features. They include cluster headache, paroxysmal hemicrania, SUNCT/SUNA, and hemicrania continua. The unifying pathophysiology involves activation of the trigeminovascular system and cranial parasympathetic outflow from the superior salivatory nucleus, with dysfunction of the hypothalamic network. Differential diagnosis among TACs is based on attack duration, frequency, periodicity, and response to indomethacin.

\subsection{Trigeminal Neuralgia}
\label{sec:Trigeminal Neuralgia}
Trigeminal neuralgia (TN, tic douloureux) is a neuropathic pain disorder characterized by brief, electric shock-like, stabbing, or shooting paroxysmal pain in one or more divisions of the trigeminal nerve (typically V2 and V3). Prevalence is 0.01--0.03\%, with a female predominance (2:1) and peak onset after age 50. The condition results from neurovascular compression of the trigeminal nerve root entry zone by a blood vessel (most commonly the superior cerebellar artery), leading to focal demyelination and ephaptic transmission; secondary TN may be caused by multiple sclerosis, tumors, or arteriovenous malformations. Clinical features include attacks lasting a fraction of a second to 2 minutes triggered by trivial stimuli (light touch, chewing, talking, toothbrushing), with no pain between attacks. Diagnosis is clinical with MRI identifying neurovascular compression. Management includes carbamazepine (200--1200 mg/day) as first-line, with a highly specific response that is itself diagnostic, or oxcarbazepine as an alternative; for refractory cases, surgical options include microvascular decompression (80--90\% long-term pain-free), stereotactic radiosurgery, or percutaneous rhizotomy. Prognosis is favorable with treatment; microvascular decompression addresses the underlying cause.

\subsection{Types of Pain}
\label{sec:Types of Pain}

Pain is classified into three mechanistic categories:
\begin{itemize}[noitemsep]
  \item \textbf{Nociceptive Pain:} Arises from activation of nociceptors by actual or threatened tissue damage. It is typically well-localized and proportional to the degree of tissue injury. Subtypes include somatic pain (sharp, aching, from skin, muscle, bone) and visceral pain (cramping, gnawing, poorly localized, from internal organs). Responds to NSAIDs and opioids.
  \item \textbf{Neuropathic Pain:} Caused by a lesion or disease of the somatosensory nervous system. Characterized by burning, shooting, electric shock-like pain, often with allodynia (pain from normally non-painful stimuli) and hyperalgesia (increased pain from normally painful stimuli). Examples include diabetic neuropathy, postherpetic neuralgia, trigeminal neuralgia. Responds to gabapentinoids, TCAs, SNRIs.
  \item \textbf{Nociplastic Pain (Central Sensitization):} Pain arising from altered nociception despite no clear evidence of actual or threatened tissue damage or disease of the somatosensory system. Defined by the IASP in 2017. Features include widespread pain, fatigue, sleep disturbance, mood disturbance, and hyperalgesia/allodynia. Central sensitization involves increased excitability of spinal and supraspinal neurons. Examples include fibromyalgia, chronic tension-type headache, and some chronic back pain. Responds to multimodal therapy including TCAs, SNRIs, exercise, and CBT.
\end{itemize}

\subsection{WHO Analgesic Ladder}
\label{sec:WHO Analgesic Ladder}

The WHO analgesic ladder, originally developed for cancer pain, is a stepwise framework for analgesic selection:
\begin{itemize}[noitemsep]
  \item \textbf{Step 1:} Non-opioid analgesics (acetaminophen, NSAIDs), with or without adjuvant analgesics, for mild pain (pain score 1--3/10).
  \item \textbf{Step 2:} Weak opioids (codeine, tramadol) or low-dose strong opioids plus non-opioids and adjuvants for moderate pain (pain score 4--6/10).
  \item \textbf{Step 3:} Strong opioids (morphine, oxycodone, hydromorphone, fentanyl, methadone) plus non-opioids and adjuvants for severe pain (pain score 7--10/10).
\end{itemize}
The ladder emphasizes regular, around-the-clock dosing with breakthrough analgesia as needed, tailored to individual patient response and side effect profile. While developed for cancer pain, it has been adapted for chronic non-cancer pain with the recognition that psychosocial factors, functional goals, and risk of opioid misuse require additional consideration.

%-------------------------------------------------------------
\section{Pain Physiology}
%-------------------------------------------------------------

%-------------------------------------------------------------

%-------------------------------------------------------------

%-------------------------------------------------------------

%-------------------------------------------------------------

%-------------------------------------------------------------

%-------------------------------------------------------------

%-------------------------------------------------------------

\label{sec:Pain Physiology}
%-------------------------------------------------------------\label{sec:Pain Physiology}

Pain is defined by the International Association for the Study of Pain (IASP) as ``an unpleasant sensory and emotional experience associated with, or resembling that associated with, actual or potential tissue damage.'' It serves as a protective mechanism but becomes pathological when it persists beyond tissue healing.

\subsection{Adjuvant Analgesics}
\label{sec:Adjuvant Analgesics}

Adjuvant analgesics are medications whose primary indication is not pain but have analgesic properties in specific pain conditions.

Gabapentinoids (gabapentin, pregabalin) bind to the alpha-2-delta subunit of voltage-gated calcium channels, reducing excitatory neurotransmitter release. They are first-line for neuropathic pain. Gabapentin is titrated from 300 mg to 900--3600 mg/day in three divided doses. Pregabalin is dosed at 150--600 mg/day in two to three divided doses. Side effects include sedation, dizziness, ataxia, peripheral edema, and weight gain. Dose adjustment is required in renal impairment.

Tricyclic antidepressants (TCAs)---amitriptyline, nortriptyline, desipramine---inhibit reuptake of serotonin and norepinephrine, and block sodium channels and NMDA receptors. They are effective for neuropathic pain, fibromyalgia, chronic tension-type headache, and migraine prevention. Doses for pain (10--150 mg at bedtime) are lower than those for depression. Side effects include anticholinergic effects (dry mouth, constipation, urinary retention, blurred vision), sedation, weight gain, and QT prolongation (caution in cardiac disease).

Serotonin-norepinephrine reuptake inhibitors (SNRIs)---duloxetine, venlafaxine, milnacipran---are effective for diabetic peripheral neuropathy, fibromyalgia, and chronic musculoskeletal pain. Duloxetine is FDA-approved for diabetic peripheral neuropathy (60--120 mg/day), fibromyalgia (60--120 mg/day), and chronic MSK pain. Venlafaxine (150--225 mg/day) has a similar profile.

Lidocaine patichtes (5\%) are topical agents effective for postherpetic neuralgia and localized neuropathic pain. The patch is applied to the painful area for 12 hours on, 12 hours off. Systemic absorption is minimal.

Capsaicin is a TRPV1 agonist that depletes substance P from peripheral nociceptors. Low-dose capsaicin creams and high-dose capsaicin 8\% patches (applied in clinic under medical supervision for 30--60 minutes) are effective for postherpetic neuralgia and peripheral neuropathic pain.

Ketamine is an NMDA receptor antagonist with potent analgesic and antihyperalgesic properties. Sub-anesthetic doses (IV infusion 0.1--0.5 mg/kg/hour) are used in acute pain (postoperative, opioid-tolerant patients, complex regional pain syndrome) and refractory chronic pain. The mechanism involves blockade of central sensitization.

\subsection{Atypical Facial Pain}
\label{sec:Atypical Facial Pain}
Persistent idiopathic facial pain (PIFP, formerly atypical facial pain) is a persistent facial pain that does not fit the diagnostic criteria for any other cranial neuralgia or headache disorder. The condition results from poorly understood pathophysiology, possibly involving central sensitization and dysfunction of descending pain modulatory systems. Clinical features include daily or nearly continuous deep, poorly localized, aching or burning pain that may cross the midline. Diagnosis requires exclusion of dental and neurological causes after thorough investigation including imaging. Management is multimodal: TCAs, SNRIs, gabapentin, pregabalin, and cognitive behavioral therapy. Prognosis is variable; invasive procedures are generally not recommended.

\subsection{Central Post-Stroke Pain}
\label{sec:Central Post-Stroke Pain}
Central post-stroke pain (CPSP, Dejerine-Roussy syndrome) is a central neuropathic pain condition arising from a cerebrovascular lesion (ischemic or hemorrhagic) affecting the spinothalamic tract or thalamocortical projections. It occurs in 1--12\% of stroke survivors, with onset typically within 1--6 months after the stroke. The condition results from the lesion in the somatosensory pathway. Clinical features include burning, aching, pricking, and freezing sensations contralateral to the lesion, with prominent allodynia and hyperalgesia to cold stimuli. Diagnosis is clinical with confirmatory neuroimaging identifying the lesion in the appropriate somatosensory pathway. First-line pharmacotherapy includes amitriptyline and lamotrigine; gabapentin, pregabalin, and duloxetine are also used. Prognosis is variable; opioids are generally ineffective for central neuropathic pain.

